\section{Further work}
\label{sec:future_work}

\subsection{Noise generalization}
\label{sec:noise_generalization}
The Gaussian perturbation is one instance of a more general conjugate
construction. Our companion manuscript \cite{conjugatecriterion2026}
characterizes symmetric scalar errors $U$ admitting an independent real
variable $V$ such that
\begin{equation}\label{eq:conjugate_transform}
M_U(t)\varphi_V(t)=1,\qquad t\in\mathbb R,
\end{equation}
where $M_U(t)=\mathbb E e^{tU}$ is finite on $\mathbb R$ and
$\varphi_V(t)=\mathbb E e^{itV}$. Existence is equivalent to $1/M_U$
being a characteristic function; the companion manuscript gives Fourier
and moment criteria and constructs the conjugate law. The imaginary
perturbation must follow this law, which generally differs from the
measurement-error law.

\paragraph{Uniform measurement error.}
For $U\sim\operatorname{Unif}[-b,b]$, with known $b>0$,
\[
M_U(t)=\frac{\sinh(bt)}{bt},\qquad
\varphi_V(t)=\frac{bt}{\sinh(bt)},
\]
with both ratios interpreted as $1$ at $t=0$.
The conjugate $V$ has a centered logistic distribution with scale $b/\pi$:
\[
f_V(v)=\frac{\pi}{4b\cosh^2(\pi v/(2b))},\qquad
F_V(v)=\frac{1}{1+e^{-\pi v/b}}.
\]
It can be sampled directly from $R\sim\operatorname{Unif}(0,1)$ via
\begin{equation}\label{eq:uniform_conjugate_sampler}
V=\frac b\pi\log\frac{R}{1-R}.
\end{equation}
Thus a uniform measurement error is corrected by an independent logistic
imaginary perturbation.

\paragraph{Symmetric two-point measurement error.}
For $\mathbb P(U=b)=\mathbb P(U=-b)=1/2$,
\[
M_U(t)=\cosh(bt),\qquad
\varphi_V(t)=\frac1{\cosh(bt)}.
\]
Here $V$ has a scaled hyperbolic-secant distribution:
\[
f_V(v)=\frac{1}{2b\cosh(\pi v/(2b))},\qquad
F_V(v)=\frac2\pi\arctan\!\left(e^{\pi v/(2b)}\right).
\]
An exact sampler is
\begin{equation}\label{eq:binary_conjugate_sampler}
V=\frac{2b}\pi\log\tan\!\left(\frac{\pi R}{2}\right),
\qquad R\sim\operatorname{Unif}(0,1).
\end{equation}
The imaginary perturbation is continuous even though the measurement
error is discrete.

\paragraph{Exact correction and implementation.}
Set $D=U+iV$. Equation~\eqref{eq:conjugate_transform} implies
$\mathbb E D^k=0$ for every integer $k\ge1$
\cite{conjugatecriterion2026}. Hence, for fixed $(\theta,z)$,
\begin{equation}\label{eq:nongaussian_correction}
\mathbb E_{U,V}\operatorname{Re}g^{\mathbb C}(\theta;z+U+iV)
=g(\theta;z)
\end{equation}
for polynomial scores. The identity also holds for entire scores when
\begin{equation}\label{eq:nongaussian_taylor}
\sum_{k=0}^\infty
\frac{\|\partial_z^k g^{\mathbb C}(\theta;z)\|}{k!}
\mathbb E|D|^k<\infty.
\end{equation}
Indeed, this condition permits termwise expectation of the Taylor series;
all nonconstant terms vanish. Thus correction is exact, with no
finite-order truncation. Algorithm~\ref{alg:nongaussian_ccsgd} implements
it using the explicit inverse-distribution samplers above. It requires
the error law and its scale, but not the latent input distribution.

\begin{algorithm}[H]
\caption{CCSGD for uniform or symmetric two-point measurement error}
\label{alg:nongaussian_ccsgd}
\begin{algorithmic}[1]
\Require Initial point $\theta_1$, step sizes $\{\eta_t\}_{t=1}^T$, known error type and scale $b>0$
\For{$t=1,\ldots,T$}
    \State Observe a fresh $W_t=Z_t+U_t$
    \State Draw $R_t\sim\operatorname{Unif}(0,1)$ independently of the observation and the past
    \If{$U_t\sim\operatorname{Unif}[-b,b]$}
        \State $V_t\gets(b/\pi)\log\{R_t/(1-R_t)\}$
    \Else\Comment{$\mathbb P(U_t=\pm b)=1/2$}
        \State $V_t\gets(2b/\pi)\log\tan(\pi R_t/2)$
    \EndIf
    \State $G_t\gets\operatorname{Re}g^{\mathbb C}(\theta_t;W_t+iV_t)$
    \State $\theta_{t+1}\gets\theta_t-\eta_tG_t$
\EndFor
\State \Return $\theta_{T+1}$
\end{algorithmic}
\end{algorithm}
For noise-free responses, the observed response is held fixed during the
complex evaluation, as in Section~\ref{sec:algorithm}. Independent
replicates of $V_t$ may be averaged using the same observation. For
independent input-error coordinates, the construction can also be applied
coordinatewise with independent conjugate perturbations, subject to the
corresponding multivariate integrability condition.

\paragraph{Moment conditions for optimization.}
Unlike a Gaussian conjugate, these conjugates have exponential moments
only below a finite threshold. Specifically,
\[
\mathbb E e^{s|V|}<\infty\quad\Longleftrightarrow\quad 0\le s<\lambda,
\qquad
\lambda=
\begin{cases}
\pi/b, & U\sim\operatorname{Unif}[-b,b],\\
\pi/(2b), & \mathbb P(U=\pm b)=1/2.
\end{cases}
\]
Under the derivative envelope~\eqref{eq:parameter_envelope}, the strict
condition $a<\lambda$ suffices for
\eqref{eq:nongaussian_taylor}: the score has exponential growth of any
rate slightly larger than $a$, and the conjugate has an exponential
moment at a still larger rate. If $2a<\lambda$, boundedness of $U$ gives
\[
K_a=\mathbb E e^{2a|Z|}\,
\mathbb E[|D|^2e^{2a|D|}]<\infty.
\]
Together with the clean sampling and amplitude bounds, the same argument
as in Section~\ref{sec:main_results} yields the relative variance
condition~\eqref{eq:relative_variance}. For the CLTs, choose $\delta>0$
with $(2+\delta)a<\lambda$; the higher-moment proof then applies as well.
Thus the optimization analysis transfers when these moment conditions
and the corresponding objective and covariance assumptions are verified.
Polynomial scores satisfy the required input moment conditions for both
pairs. Arbitrary entire scores still require growth control, and random
responses require the corresponding joint moment bounds.

\subsection{Function generalization}
Exact Gaussian correction restricts the admissible score class. In the scalar Gaussian location model with positive noise
variance, Stefanski \cite{stefanski1989} shows that a target admitting an
integrable unbiased estimator for every real mean must extend to an entire
function. Applied componentwise at fixed $\theta$, this rules out globally
exact unbiased correction of non-entire scores in that setting.
Entire analyticity alone does not ensure the integrability required by
CCSGD; growth control remains essential. For non-entire scores, a natural
direction is approximate or regularized correction with explicit bias,
variance, and optimization-error bounds. With singular $\Omega$, the
analyticity restriction concerns the directions affected by noise.
